\documentclass[11pt,a4paper]{article}
\usepackage[margin=2.5cm]{geometry}
\usepackage{amsmath,amssymb}
\usepackage{booktabs}
\usepackage{enumitem}
\usepackage{hyperref}
\usepackage{xcolor}

\title{Spatial Memory Field}
\author{Social Dynamics Simulation}
\date{}

\begin{document}
\maketitle

\section{Notation}

\begin{tabular}{ll}
\toprule
\textbf{Symbol} & \textbf{Definition} \\
\midrule
$N$ & Number of particles \\
$K$ & Number of preference dimensions per particle \\
$G$ & Grid resolution ($G \times G$ cells covering the domain) \\
$L$ & Domain size (toroidal, default $L = 1.0$) \\
$i$ & Particle index, $i \in \{1, \ldots, N\}$ \\
$d$ & Preference dimension index, $d \in \{1, \ldots, K\}$ \\
$\mathbf{x}_i \in [0, L)^2$ & Position of particle $i$ \\
$\mathbf{p}_i \in [-1, 1]^K$ & Preference (signal) vector of particle $i$ \\
$\tilde{\mathbf{p}}_i$ & Effective preference after field modulation \\
$M(c, d)$ & Memory field value at grid cell $c$, dimension $d$ \\
$\mathrm{cell}(\mathbf{x})$ & Grid cell containing position $\mathbf{x}$: $\left(\lfloor \mathbf{x} \cdot G/L \rfloor\right) \bmod G$ \\
$\alpha$ & Field strength (how strongly the field modulates preferences) \\
$\beta$ & Write rate (how fast particles deposit into the field) \\
$\gamma$ & Decay factor per step ($0 < \gamma \leq 1$) \\
$\sigma$ & Gaussian blur standard deviation (in grid cells) \\
$s$ & Social learning rate (positive = conformity, negative = differentiation) \\
\bottomrule
\end{tabular}

\section{Overview}

A persistent grid field $M(\mathbf{x}, d)$ of shape $(G \times G \times K)$ creates a feedback loop between particle dynamics and the environment. Particles both \textbf{read} from and \textbf{write} to the field as they move through space, allowing the system to develop location-dependent behavioral biases based on interaction history.

\section{Mechanism}

\subsection{Step 1: Read (before physics)}

Each particle $i$'s preferences are modulated by the field at its grid cell:
\begin{equation}
\tilde{p}_i^{(d)} = p_i^{(d)} \cdot \left(1 + \alpha \cdot M\!\left(\mathrm{cell}(\mathbf{x}_i),\, d\right)\right)
\end{equation}

where:
\begin{itemize}[nosep]
\item $p_i^{(d)}$ is particle $i$'s preference in dimension $d$
\item $\tilde{p}_i^{(d)}$ is the effective (modulated) preference used for physics
\item $\alpha$ is the \textbf{field strength} parameter
\item $M\!\left(\mathrm{cell}(\mathbf{x}_i), d\right)$ is the memory field value at particle $i$'s grid cell
\end{itemize}

The modulated preferences are clipped to $[-1, 1]$.

\textbf{Effect}: Positive field values amplify the matching preference dimension. Negative values dampen it. A region where $M(\mathbf{x}, 0) > 0$ will cause particles with $p^{(0)} > 0$ to have stronger attraction in dimension~0, while particles with $p^{(0)} < 0$ will have weaker (dampened) repulsion.

\subsection{Step 2: Physics}

The simulation runs normally using the modulated preferences $\tilde{\mathbf{p}}_i$ for neighbor selection and compatibility weighting. The original (unmodulated) preferences $\mathbf{p}_i$ are preserved.

Any social learning delta $\Delta \mathbf{p}$ from the physics kernel is captured and applied to the \textbf{original} preferences:
\begin{equation}
\mathbf{p}_i \leftarrow \mathrm{clip}\!\left(\mathbf{p}_i + \Delta \mathbf{p},\; -1,\; 1\right)
\end{equation}

This ensures the field only affects interaction forces, not the stored identity of the particles.

\subsection{Step 3: Write (after physics)}

Each particle $i$ deposits its raw preferences into the field at its \textbf{new} position:
\begin{equation}
M\!\left(\mathrm{cell}(\mathbf{x}_i),\, d\right) \;\leftarrow\; M\!\left(\mathrm{cell}(\mathbf{x}_i),\, d\right) + \beta \cdot p_i^{(d)}
\end{equation}

where $\beta$ is the \textbf{write rate}. Multiple particles in the same cell accumulate additively.

\subsection{Step 4: Decay}

The entire field decays exponentially:
\begin{equation}
M \leftarrow \gamma \cdot M
\end{equation}

where $\gamma$ is the \textbf{decay} parameter. The half-life in steps is:
\begin{equation}
t_{1/2} = \frac{\ln 2}{\ln(1/\gamma)}
\end{equation}

\subsection{Step 5: Blur (optional)}

A Gaussian blur with standard deviation $\sigma$ (in grid cells) is applied to each dimension independently, with periodic boundary wrapping:
\begin{equation}
M^{(d)} \leftarrow M^{(d)} * \mathcal{G}(\sigma)
\end{equation}

This diffuses the memory spatially, causing particle traces to spread into neighboring cells.

\section{Parameters}

\begin{tabular}{lcccp{6cm}}
\toprule
\textbf{Parameter} & \textbf{Symbol} & \textbf{Range} & \textbf{Default} & \textbf{Description} \\
\midrule
Write Rate & $\beta$ & $[0, 1]$ & 0.01 & Deposition rate per particle per step \\
Field Strength & $\alpha$ & $[0, 5]$ & 0.5 & Modulation strength on preferences \\
Field Decay & $\gamma$ & $[0.9, 1]$ & 0.999 & Exponential decay per step \\
Blur Sigma & $\sigma$ & $[0.1, 10]$ & 1.0 & Gaussian diffusion radius (grid cells) \\
\bottomrule
\end{tabular}

\subsection{Decay half-lives}

\begin{tabular}{crl}
\toprule
$\gamma$ & \textbf{Half-life (steps)} & \textbf{Character} \\
\midrule
0.99 & 69 & Short-term memory \\
0.995 & 139 & Medium-term \\
0.999 & 693 & Long-term memory \\
0.9999 & 6931 & Near-permanent \\
1.0 & $\infty$ & Permanent (grows without bound) \\
\bottomrule
\end{tabular}

\section{Expected Behaviors}

\subsection{Regime 1: Faint trails (low $\beta$, low $\alpha$)}

Particles leave barely perceptible traces. The Memory view shows thin paths where particles have traveled. Minimal feedback on dynamics---the system behaves almost identically to no memory.

\textbf{Settings}: $\beta = 0.001$, $\alpha = 0.1$, $\gamma = 0.999$, blur off.

\subsection{Regime 2: Territory formation (high $\beta$, high $\alpha$, no blur)}

Rapid spatial polarization. The first particles to visit a region ``claim'' it:
\begin{enumerate}[nosep]
\item Particle with $p^{(0)} = +1$ passes through cell $(\mathbf{x})$
\item Field accumulates $M(\mathbf{x}, 0) > 0$
\item Future particles at $\mathbf{x}$ get $\tilde{p}^{(0)}$ amplified if positive, dampened if negative
\item Positive-$p^{(0)}$ particles are attracted more strongly $\rightarrow$ they stay $\rightarrow$ more deposition
\end{enumerate}

This positive feedback creates sharp boundaries between regions of different preference types.

\textbf{Settings}: $\beta = 0.05$, $\alpha = 1.0$, $\gamma = 0.999$, blur off.

\subsection{Regime 3: Pheromone landscapes (medium $\beta$, medium $\alpha$, with blur)}

Smooth preference gradients emerge instead of sharp territories. The blur diffuses deposits into broad regional biases. Particles drift along these gradients, reinforcing them.

This is closest to \textbf{ant colony pheromone dynamics}---particles leave chemical-like traces that guide future behavior, with spatial diffusion and temporal decay.

\textbf{Settings}: $\beta = 0.01$, $\alpha = 0.5$, $\gamma = 0.999$, blur on, $\sigma = 2.0$.

\subsection{Regime 4: Memory + differentiation (high $\alpha$ + negative social)}

The most interesting regime. \textbf{Social learning} ($s < 0$) pushes particles to differentiate from their neighbors, while the \textbf{memory field} pushes them to conform to local history.

\begin{itemize}[nosep]
\item Social force: ``be different from who's near you now''
\item Memory force: ``be similar to who was here before''
\end{itemize}

The tension between these opposing pressures can create:
\begin{itemize}[nosep]
\item \textbf{Oscillations}: particles differentiate, leave, field decays, new particles arrive and conform to old memory, then differentiate again
\item \textbf{Traveling waves}: regions of preference bias propagate as particles carry information
\item \textbf{Persistent structures}: stable spatial patterns where the field and particle distribution reach equilibrium
\end{itemize}

\textbf{Settings}: $\beta = 0.01$, $\alpha = 1.0$, $\gamma = 0.999$, $s = -0.003$, blur on, $\sigma = 1.0$.

\subsection{Regime 5: Colony ghosts (high $\alpha$, medium $\gamma$)}

When particles cluster and deposit strongly, the field at those cells grows, amplifying preferences, increasing attraction, tightening the cluster. This is a \textbf{positive feedback loop} that creates stable ``colonies.''

If the cluster eventually disperses (due to repulsion or social differentiation), the field retains a ``ghost'' of the colony---a region of strong field values with no particles. This ghost can:
\begin{itemize}[nosep]
\item Attract future particles of the same type back to the abandoned site
\item Create ``territorial memory'' that persists after the original occupants leave
\item Generate cyclic colonization--abandonment patterns
\end{itemize}

\textbf{Settings}: $\beta = 0.02$, $\alpha = 1.5$, $\gamma = 0.998$, blur on, $\sigma = 1.0$.

\section{Visualization}

The \textbf{Memory} right-panel view shows $M(\mathbf{x}, d)$ as RGB:
\begin{itemize}[nosep]
\item Red channel: $M(\mathbf{x}, 0)$
\item Green channel: $M(\mathbf{x}, 1)$
\item Blue channel: $M(\mathbf{x}, 2)$
\end{itemize}

Each channel is independently normalized by its maximum absolute value, mapped from $[-|M|_{\max}, +|M|_{\max}]$ to $[0, 1]$:
\begin{itemize}[nosep]
\item Gray (0.5, 0.5, 0.5) = neutral (no memory)
\item Bright color = strong positive field in that dimension
\item Dark color = strong negative field
\item White = all dimensions strongly positive
\end{itemize}

\section{Mathematical Properties}

\subsection{Fixed point analysis}

At steady state ($\dot{M} = 0$), balancing deposition against decay gives:
\begin{equation}
M^* = \frac{\beta}{1 - \gamma} \cdot \langle p \rangle_{\mathrm{local}}
\end{equation}

where $\langle p \rangle_{\mathrm{local}}$ is the time-averaged preference of particles passing through each cell. For $\gamma = 0.999$ and $\beta = 0.01$:
\begin{equation}
M^* \approx \frac{0.01}{0.001} \cdot \langle p \rangle_{\mathrm{local}} = 10 \cdot \langle p \rangle_{\mathrm{local}}
\end{equation}

So the steady-state field is approximately $10\times$ the local average preference.

\subsection{Feedback gain}

The effective feedback gain is $\alpha \cdot |M^*|$. For the modulation to be significant but not overwhelming:
\begin{equation}
\alpha \cdot |M^*| \sim \mathcal{O}(1)
\end{equation}

With $M^* \approx 10 \cdot \langle p \rangle$ and $|\langle p \rangle| \sim 0.5$:
\begin{equation}
\alpha \cdot 5 \sim 1 \quad\implies\quad \alpha \sim 0.2
\end{equation}

This suggests $\alpha \approx 0.2$ is the transition point between weak and strong feedback.

\subsection{Diffusion length}

With blur sigma $\sigma$ applied every step with decay $\gamma$, the effective diffusion length (how far a deposit spreads before decaying to $1/e$) is:
\begin{equation}
L_{\mathrm{diff}} = \sigma \cdot \sqrt{\frac{2}{1 - \gamma}} \cdot \frac{L}{G}
\end{equation}

For $\sigma = 2$, $\gamma = 0.999$, $G = 256$:
\begin{equation}
L_{\mathrm{diff}} \approx 2 \cdot \sqrt{2000} \cdot \frac{1}{256} \approx 0.35
\end{equation}

This means a deposit diffuses across about 35\% of the domain before decaying. Higher $\sigma$ or higher $\gamma$ extends this.

\end{document}
